\documentclass{article}
\usepackage{amsmath,amssymb,amsthm}
\usepackage{algorithm,algorithmic}
\usepackage{enumitem}
\usepackage{hyperref}
\newtheorem{theorem}{Theorem}
\newtheorem{lemma}{Lemma}
\newtheorem{assumption}{Assumption}
\newcommand{\norm}[1]{\left\lVert#1\right\rVert}
\newcommand{\iprod}[2]{\left\langle #1,#2\right\rangle}
\begin{document}

\section*{Primal–Dual Hybrid Gradient (PDHG)}

\section*{Mathematical Formulation}
Consider the convex–concave saddle‑point problem  

\[
\min_{x\in\mathbb R^{n}} \max_{y\in\mathbb R^{m}} \; \mathcal L(x,y)
:= f(x)+\langle Kx ,y\rangle - g^{*}(y),
\tag{1}
\]

where  

\begin{itemize}[nosep]
\item $f:\mathbb R^{n}\!\to\!\mathbb R$ is convex and $L$–smooth,
\item $g^{*}:\mathbb R^{m}\!\to\!\mathbb R\cup\{+\infty\}$ is convex (the Fenchel conjugate of a proper closed convex $g$),
\item $K\in\mathbb R^{m\times n}$ is a linear operator with operator norm $\|K\|$.
\end{itemize}

The PDHG iteration with primal step size $\tau>0$ and dual step size $\sigma>0$ reads  

\[
\boxed{
\begin{aligned}
y^{k+1} &= \operatorname{prox}_{\sigma g^{*}}\!\Bigl(y^{k}+ \sigma K x^{k}\Bigr), \tag{2a}\\[4pt]
x^{k+1} &= \operatorname{prox}_{\tau f}\!\Bigl(x^{k}-\tau K^{\!\top} y^{k+1}\Bigr). \tag{2b}
\end{aligned}}
\]

The proximal operators are defined by  

\[
\operatorname{prox}_{\lambda h}(u):=
\arg\min_{z}\Bigl\{h(z)+\frac{1}{2\lambda}\|z-u\|^{2}\Bigr\}.
\tag{3}
\]

When $f$ is $L$‑smooth, $\operatorname{prox}_{\tau f}$ admits the explicit form  

\[
\operatorname{prox}_{\tau f}(u)=
\bigl(I+\tau \nabla^{2}f\bigr)^{-1}\!\bigl(u-\tau \nabla f(u)\bigr),
\tag{4}
\]

and for $g^{*}=0$ (i.e. $g\equiv 0$) the dual proximal step reduces to a simple gradient ascent:  

\[
y^{k+1}=y^{k}+ \sigma K x^{k}. \tag{5}
\]

\section*{Pseudocode}
\begin{algorithm}[H]
\caption{Primal–Dual Hybrid Gradient}
\begin{algorithmic}[1]
\STATE \textbf{Input:} $x^{0}\in\mathbb R^{n}$, $y^{0}\in\mathbb R^{m}$, stepsizes $\tau,\sigma>0$.
\FOR{$k=0,1,2,\dots$}
\STATE $y^{k+1}\gets \operatorname{prox}_{\sigma g^{*}}\!\bigl(y^{k}+ \sigma K x^{k}\bigr)$ \hfill\textit{(dual ascent)}
\STATE $x^{k+1}\gets \operatorname{prox}_{\tau f}\!\bigl(x^{k}-\tau K^{\!\top} y^{k+1}\bigr)$ \hfill\textit{(primal descent)}
\ENDFOR
\STATE \textbf{Output:} last iterates $(x^{k},y^{k})$ or ergodic averages.
\end{algorithmic}
\end{algorithm}

\section*{Dry Run Example}
We solve the scalar problem  

\[
\min_{x\in\mathbb R}\;(x-3)^{2},
\]

which fits (1) with  

\[
f(x)=(x-3)^{2},\qquad g^{*}\equiv0,\qquad K=1.
\]

The gradient $\nabla f(x)=2(x-3)$, $L=2$.
Choose $\tau=\sigma=0.1$ (satisfying $\tau\sigma\|K\|^{2}=0.01<1$) and start from $x^{0}=0$, $y^{0}=0$.

\begin{center}
\begin{tabular}{c|c|c|c}
$k$ & $y^{k+1}=y^{k}+ \sigma x^{k}$ & $x^{k+1}= \operatorname{prox}_{\tau f}(x^{k}-\tau y^{k+1})$ & $f(x^{k+1})$\\ \hline
0 & $y^{1}=0+0.1\cdot0=0$ &
\begin{aligned}
u^{0}&=x^{0}-\tau y^{1}=0,\\
x^{1}&=\frac{u^{0}+2\tau\cdot3}{1+2\tau}
      =\frac{0+0.6}{1.2}=0.5
\end{aligned}
& $(0.5-3)^{2}=6.25$\\ \hline
1 & $y^{2}=0+0.1\cdot0.5=0.05$ &
\begin{aligned}
u^{1}&=x^{1}-\tau y^{2}=0.5-0.005=0.495,\\
x^{2}&=\frac{u^{1}+0.6}{1.2}= \frac{1.095}{1.2}=0.9125
\end{aligned}
& $(0.9125-3)^{2}\approx4.36$\\ \hline
2 & $y^{3}=0.05+0.1\cdot0.9125=0.14125$ &
\begin{aligned}
u^{2}&=x^{2}-\tau y^{3}=0.9125-0.014125=0.898375,\\
x^{3}&=\frac{u^{2}+0.6}{1.2}= \frac{1.498375}{1.2}=1.248646
\end{aligned}
& $(1.248646-3)^{2}\approx3.07$
\end{tabular}
\end{center}

The objective value decreases, illustrating the algorithm’s progress.

\section*{Assumptions}
\begin{assumption}[Convexity and Smoothness]\label{as:convex}
$f:\mathbb R^{n}\!\to\!\mathbb R$ is convex and differentiable with $L$‑Lipschitz gradient:
\[
\|\nabla f(x)-\nabla f(x')\|\le L\|x-x'\|,\qquad\forall x,x'.
\]
$g^{*}$ is convex, proper, and lower‑semicontinuous.
\end{assumption}

\begin{assumption}[Step‑size Condition]\label{as:steps}
The primal and dual steps satisfy  
\[
\tau\sigma\|K\|^{2}<1.
\tag{6}
\]
\end{assumption}

\begin{assumption}[Existence of Saddle Point]\label{as:saddle}
There exists $(x^{\star},y^{\star})$ such that  
\[
\mathcal L(x^{\star},y)\le \mathcal L(x^{\star},y^{\star})\le \mathcal L(x,y^{\star}),\qquad\forall x,y.
\]
\end{assumption}

\section*{Main Convergence Theorem}
\begin{theorem}[Ergodic $O(1/k)$ Convergence]\label{thm:ergodic}
Let Assumptions \ref{as:convex}–\ref{as:saddle} hold and choose $\tau,\sigma$ obeying \eqref{as:steps}. Define the ergodic averages  

\[
\bar x^{K}:=\frac{1}{K}\sum_{k=0}^{K-1}x^{k},\qquad 
\bar y^{K}:=\frac{1}{K}\sum_{k=0}^{K-1}y^{k}.
\]

Then for any $(x,y)$ we have  

\[
\underbrace{\mathcal L(\bar x^{K},y)-\mathcal L(x,\bar y^{K})}_{\text{primal–dual gap}}
\le
\frac{1}{2K}\Bigl(\frac{\|x^{0}-x\|^{2}}{\tau}
      +\frac{\|y^{0}-y\|^{2}}{\sigma}\Bigr).
\tag{7}
\]

In particular, taking $(x,y)=(x^{\star},y^{\star})$ yields  

\[
\mathcal L(\bar x^{K},y^{\star})-\mathcal L(x^{\star},\bar y^{K})
\le \frac{C}{K},\qquad C:=\frac{1}{2}\Bigl(\frac{\|x^{0}-x^{\star}\|^{2}}{\tau}
      +\frac{\|y^{0}-y^{\star}\|^{2}}{\sigma}\Bigr).
\tag{8}
\]
Thus the primal–dual gap decays as $O(1/K)$.
\end{theorem}

\section*{Theoretical Properties}
\begin{enumerate}[nosep]
\item \textbf{Stability.} Condition \eqref{as:steps} guarantees the Lyapunov function  

\[
\Phi^{k}:=\frac{1}{2\tau}\|x^{k}-x^{\star}\|^{2}
           +\frac{1}{2\sigma}\|y^{k}-y^{\star}\|^{2}
\]

is non‑increasing up to the gap term (see Lemma~\ref{lem:descent}).
\item \textbf{Hyper‑parameter Sensitivity.} The bound \eqref{7} depends inversely on $\tau$ and $\sigma$; larger steps accelerate the constant $C$ but must respect $\tau\sigma\|K\|^{2}<1$.
\item \textbf{Comparison to Baselines.}
\begin{itemize}[nosep]
\item Compared with vanilla Gradient Descent on the primal alone, PDHG attains the same $O(1/k)$ rate while handling the coupling term $\langle Kx,y\rangle$ without explicit matrix inversion.
\item When $g^{*}=0$ and $K=I$, PDHG reduces to the celebrated Chambolle–Pock scheme, whose $O(1/k)$ rate is known to be optimal for general convex–concave saddle problems.
\end{itemize}
\end{enumerate}

\section*{Discussion}
The theorem shows that PDHG converges in the ergodic sense without requiring strong convexity. The linear rate $O(1/k)$ matches the optimal rate for first‑order methods on convex–concave saddle points. Extensions (e.g., over‑relaxation, acceleration) can improve constants but retain the same asymptotic order.

\section*{Preliminaries (Definitions and Lemmas)}
\begin{lemma}[Three‑Point Identity for Prox]\label{lem:prox}
For any closed convex $h$, any $\lambda>0$, and any $u,v$,
\[
\langle \operatorname{prox}_{\lambda h}(u)-u,\;z-\operatorname{prox}_{\lambda h}(u)\rangle
\ge
\lambda\bigl(h(z)-h(\operatorname{prox}_{\lambda h}(u))\bigr),\qquad\forall z.
\tag{9}
\]
\end{lemma}
\begin{proof}
Standard sub‑gradient optimality condition of the proximal minimization; see e.g. \cite{Bauschke2011}.
\end{proof}

\begin{lemma}[Descent Lemma for $L$‑smooth $f$]\label{lem:descent}
For any $x,x'$,
\[
f(x')\le f(x)+\langle \nabla f(x),x'-x\rangle+\frac{L}{2}\|x'-x\|^{2}.
\tag{10}
\]
\end{lemma}
\begin{proof}
Direct consequence of $L$‑Lipschitz gradient.
\end{proof}

\section*{Main Proof (Step‑by‑Step Derivation)}
We prove Theorem~\ref{thm:ergodic}. Throughout the proof we fix an arbitrary saddle point $(x^{\star},y^{\star})$.

\begin{enumerate}
\item \textbf{Dual proximal inequality.}  
Apply Lemma~\ref{lem:prox} with $h=g^{*}$, $\lambda=\sigma$, $u=y^{k}+ \sigma Kx^{k}$ and $z=y^{\star}$:
\[
\langle y^{k+1}-y^{k}-\sigma Kx^{k},\; y^{\star}-y^{k+1}\rangle
\ge \sigma\bigl(g^{*}(y^{\star})-g^{*}(y^{k+1})\bigr).
\tag{Step 1}
\]
Rearranging,
\[
\langle y^{k+1}-y^{k},y^{\star}-y^{k+1}\rangle
\ge \sigma\bigl(g^{*}(y^{\star})-g^{*}(y^{k+1})\bigr)
+\sigma\langle Kx^{k},y^{k+1}-y^{\star}\rangle .
\tag{Step 2}
\]

\item \textbf{Primal proximal inequality.}  
Apply Lemma~\ref{lem:prox} with $h=f$, $\lambda=\tau$, $u=x^{k}-\tau K^{\!\top}y^{k+1}$ and $z=x^{\star}$:
\[
\langle x^{k+1}-x^{k}+\tau K^{\!\top}y^{k+1},\; x^{\star}-x^{k+1}\rangle
\ge \tau\bigl(f(x^{\star})-f(x^{k+1})\bigr).
\tag{Step 3}
\]
Equivalently,
\[
\langle x^{k+1}-x^{k},x^{\star}-x^{k+1}\rangle
\ge \tau\bigl(f(x^{\star})-f(x^{k+1})\bigr)
-\tau\langle K^{\!\top}y^{k+1},x^{k+1}-x^{\star}\rangle .
\tag{Step 4}
\]

\item \textbf{Combine the two inner products.}  
Add \textbf{Step 2} and \textbf{Step 4} and use the identity  
$\langle a,b\rangle =\frac12\bigl(\|a\|^{2}+\|b\|^{2}-\|a-b\|^{2}\bigr)$ to rewrite each term:

\[
\begin{aligned}
&\frac{1}{2\tau}\bigl(\|x^{k}-x^{\star}\|^{2}-\|x^{k+1}-x^{\star}\|^{2}
-\|x^{k+1}-x^{k}\|^{2}\bigr)\\
&\quad+\frac{1}{2\sigma}\bigl(\|y^{k}-y^{\star}\|^{2}-\|y^{k+1}-y^{\star}\|^{2}
-\|y^{k+1}-y^{k}\|^{2}\bigr)\\
&\quad+\langle Kx^{k},y^{k+1}-y^{\star}\rangle
-\langle K^{\!\top}y^{k+1},x^{k+1}-x^{\star}\rangle
\\
&\ge f(x^{\star})-f(x^{k+1})+g^{*}(y^{\star})-g^{*}(y^{k+1}).
\end{aligned}
\tag{Step 5}
\]

\item \textbf{Coupling term simplification.}  
Observe that  

\[
\langle Kx^{k},y^{k+1}-y^{\star}\rangle
-\langle K^{\!\top}y^{k+1},x^{k+1}-x^{\star}\rangle
=
\langle K(x^{k}-x^{k+1}),y^{k+1}-y^{\star}\rangle .
\tag{Step 6}
\]

Apply Cauchy–Schwarz and Young’s inequality with a parameter $\theta>0$ to bound the right‑hand side:

\[
\langle K(x^{k}-x^{k+1}),y^{k+1}-y^{\star}\rangle
\le \frac{\theta}{2}\|K(x^{k}-x^{k+1})\|^{2}
   +\frac{1}{2\theta}\|y^{k+1}-y^{\star}\|^{2}.
\tag{Step 7}
\]

Using $\|Kz\|\le\|K\|\|z\|$ and choosing $\theta=\tau\sigma\|K\|^{2}$ (which is $<1$ by Assumption~\ref{as:steps}), we obtain

\[
\langle K(x^{k}-x^{k+1}),y^{k+1}-y^{\star}\rangle
\le \frac{\tau\sigma\|K\|^{2}}{2}\|x^{k+1}-x^{k}\|^{2}
   +\frac{1}{2\tau\sigma\|K\|^{2}}\|y^{k+1}-y^{\star}\|^{2}.
\tag{Step 8}
\]

\item \textbf{Insert the bound into \textbf{Step 5}.}  
After moving the two norm terms to the left‑hand side, we get

\[
\begin{aligned}
&\frac{1}{2\tau}\bigl(\|x^{k}-x^{\star}\|^{2}
   -\|x^{k+1}-x^{\star}\|^{2}\bigr)
 +\frac{1}{2\sigma}\bigl(\|y^{k}-y^{\star}\|^{2}
   -\|y^{k+1}-y^{\star}\|^{2}\bigr)\\
&\qquad\le f(x^{k+1})-f(x^{\star})+g^{*}(y^{k+1})-g^{*}(y^{\star})
   +\underbrace{\Bigl(\frac{1}{2}-\frac{\tau\sigma\|K\|^{2}}{2}\Bigr)}_{\ge 0}
      \|x^{k+1}-x^{k}\|^{2}
   +\underbrace{\Bigl(\frac{1}{2}-\frac{1}{2\tau\sigma\|K\|^{2}}\Bigr)}_{<0}
      \|y^{k+1}-y^{\star}\|^{2}.
\end{aligned}
\tag{Step 9}
\]

Because $\tau\sigma\|K\|^{2}<1$, the coefficient in front of $\|y^{k+1}-y^{\star}\|^{2}$ is non‑positive and can be dropped, yielding the clean inequality

\[
\frac{1}{2\tau}\bigl(\|x^{k}-x^{\star}\|^{2}
   -\|x^{k+1}-x^{\star}\|^{2}\bigr)
 +\frac{1}{2\sigma}\bigl(\|y^{k}-y^{\star}\|^{2}
   -\|y^{k+1}-y^{\star}\|^{2}\bigr)
\le \mathcal L(x^{k+1},y^{\star})-\mathcal L(x^{\star},y^{k+1}).
\tag{Step 10}
\]

\item \textbf{Summation over iterations.}  
Sum \textbf{Step 10} for $k=0,\dots,K-1$:

\[
\frac{1}{2\tau}\|x^{0}-x^{\star}\|^{2}
+\frac{1}{2\sigma}\|y^{0}-y^{\star}\|^{2}
\ge \sum_{k=0}^{K-1}\bigl[\mathcal L(x^{k+1},y^{\star})-\mathcal L(x^{\star},y^{k+1})\bigr].
\tag{Step 11}
\]

\item \textbf{Use convex–concave structure.}  
By convexity in $x$ and concavity in $y$,
\[
\mathcal L(\bar x^{K},y^{\star})\le\frac{1}{K}\sum_{k=1}^{K}\mathcal L(x^{k},y^{\star}),\qquad
\mathcal L(x^{\star},\bar y^{K})\ge\frac{1}{K}\sum_{k=1}^{K}\mathcal L(x^{\star},y^{k}).
\tag{Step 12}
\]

Combining \textbf{Step 11} and \textbf{Step 12} gives

\[
\mathcal L(\bar x^{K},y^{\star})-\mathcal L(x^{\star},\bar y^{K})
\le\frac{1}{2K}\Bigl(\frac{\|x^{0}-x^{\star}\|^{2}}{\tau}
+\frac{\|y^{0}-y^{\star}\|^{2}}{\sigma}\Bigr),
\tag{Step 13}
\]
which is exactly the bound \eqref{7} claimed in Theorem~\ref{thm:ergodic}.
\end{enumerate}

Thus the theorem is proved. \hfill $\square$

\section*{Rate Derivation (With Constants)}
From \eqref{Step 13} define  

\[
C:=\frac12\Bigl(\frac{\|x^{0}-x^{\star}\|^{2}}{\tau}
+\frac{\|y^{0}-y^{\star}\|^{2}}{\sigma}\Bigr).
\]

Then for every $K\ge1$  

\[
\mathcal L(\bar x^{K},y^{\star})-\mathcal L(x^{\star},\bar y^{K})\le \frac{C}{K}.
\tag{14}
\]

If additionally $f$ is $\mu$‑strongly convex, the distance $\| \bar x^{K}-x^{\star}\|$ inherits an $O(1/K)$ rate, but the primal–dual gap already enjoys the optimal $O(1/K)$ without strong convexity.

\section*{Special Cases}
\begin{enumerate}[nosep]
\item \textbf{Pure primal gradient descent.}  
When $g^{*}=0$ and $K=0$, the dual updates are inert ( $y^{k}=y^{0}$ ) and \eqref{2b} reduces to $x^{k+1}= \operatorname{prox}_{\tau f}(x^{k})$, i.e. standard gradient descent with step $\tau$.
\item \textbf{Chambolle–Pock scheme.}  
For $f$ and $g$ proper convex and $K$ arbitrary, PDHG coincides with the Chambolle–Pock algorithm. The same proof yields the classical $O(1/k)$ ergodic rate.
\item \textbf{Accelerated PDHG.}  
If one introduces an over‑relaxation parameter $\theta\in(0,1]$ and replaces $x^{k}$ by $\tilde x^{k}=x^{k}+\theta (x^{k}-x^{k-1})$ in the dual step, the proof can be adapted to obtain the accelerated rate $O(1/k^{2})$ under additional smoothness of $g^{*}$.
\end{enumerate}

\begin{thebibliography}{9}
\bibitem{Bauschke2011}
H.~H. Bauschke and P.~L. Combettes,
\textit{Convex Analysis and Monotone Operator Theory in Hilbert Spaces},
Springer, 2011.
\end{thebibliography}

\end{document}